\documentclass[a4paper,12pt,oneside,pointlessnumbers]{scrartcl} 
\input{preambel} 	% Alle Packages und Formatierungen werden in preamble.tex eingetragen um den Text hier übersichtlicher zu gestalten.
% * <chemie.woski@gmail.com> 2016-05-21T11:11:27.143Z:

\begin{document}
 
%\input{title}			%Schöneres Deckblatt ohne Note
\input{title2}			%leicht modifiziertes offizielles Deckblatt vom KM
 
\begingroup
\pagestyle{empty} % ohne Seitenzahl
\tableofcontents % Inhaltsverzeichnis
\clearpage % Seite bis zum Ende leeren
\endgroup

\newpage

\setcounter{page}{1} % Seitenzähler wird auf 0 gesetzt

\part{Einleitung} % Überschrift für erstes Kapitel, optional auch: [Eintrag im IHV]{Einleitung}
\label{sec:Einleitung} % Textmarke, um darauf zu verweisen.
\lipsum[2]			% Dummytext (einfach löschen und los schreiben)

\part{Hauptteil}
\section{Thematik in der Literatur}
\lipsum[30]				% Dummytext (einfach löschen und los schreiben)

Es gibt eine Reihe von Möglichkeiten chemische Formeln wie in \textsl{\ref{XXX}} zu notieren. Eine davon ist diese \cite[S.\,122]{demo}: 

Dies ist ein Absatz und dahinter kommt der Literaturhinweis.\citeup{demo, demo2,demo3}{}

\begin{equation}					% Gleichung
\begin{aligned}						% Nummerierung am Rand
\label{XXX}							% beliebiges Label
\ce{XX2 + Y ->[{dr"uber}][{drunter}] X2 + XY}
\end{aligned}
\end{equation}
\vspace{0.2cm}

Wenn nur einzelne Ionen notiert werden sollen können diese entweder direkt im Text als $\anion{--II}{O}{}{2}$ geschrieben, oder in der Reaktionsgleichung \textsl{\ref{YYY}} notiert werden ... in Literatur \cite{latexit}:

\begin{equation}					% Gleichung
\begin{aligned}						% Nummerierung am Rand
\label{YYY}							% beliebiges Label
\textnormal{O}_2 + 4\, \ee{}{}\ch{->[{Reduktion}][{}]} \hspace{0.3cm} 2\,\anion{--II}{O}{}{2} 
\end{aligned}
\end{equation}
\vspace{0.2cm}

\begin{figure}[!ht]		% Bild mit Unterschrift einbinden
\centering
\includegraphics[scale=1]{pics/bild.png} 
\caption[Demobild]{Einzelne Bilder wie dieses sollten zentriert auf der Seite platziert werden. Auf eine sinnvolle Größe ist ebenso zu achten, wie auf eine konkrete Beschreibung der Abbildung.}
\label{fig:bild1}\vspace{0.8em}
\end{figure}

\lipsum[31]
Verweis auf das Bild \ref{fig:bild1}.

\newpage
\section{Überschrift}

\renewcommand{\baselinestretch}{1.0}\normalsize 
\setlength{\columnsep}{30pt} % Abstand Spalten vergrößern
\begin{multicols}{2}[\noindent \textbf{Hier könnte beispielsweise ein Verweis auf den Textauszug aus der Literatur stehen.\citeup{voss}{}}] %\textbf{Dies kann als Zusammenfassung für den darunter liegenden Text gesehen werden. Hin und wieder ist dieser auch fett gedruckt.}
\setcounter{linenumber}{1}
\begin{linenumbers} % Schaltet Zeilennummerierung ein
%\modulolinenumbers[5] % nur jede 5. Zeile
\lipsum[99]
\end{linenumbers}
\end{multicols} % Ende 2 Spalten
\renewcommand{\baselinestretch}{1.5}\normalsize 

\subsection{Überschrift}
\lipsum[99]



\lipsum[99]

\newpage
\part{Schluss}
Dies ist eine Text\footnote{Dies solle Fußnote werden, die den Rahmen einer Zeile sprengt. Dies lässt sich auch durch Wiederholung den Inhalts erreichen. Dies lässt sich auch durch Wiederholung den Inhalts erreichen.  Dies lässt sich auch durch Wiederholung den Inhalts erreichen. \newline Dies lässt sich auch durch Wiederholung den Inhalts erreichen.}, der auch in der Literatur existiert.\citeup{demo}{,S.22} 
\lipsum[32]

Da pack ich doch gleich no ne Fußnote drunter.\footnote{Mal sehen, welcher Abstand hier zu vorherigen Fußnote entsteht.}

\lipsum[10]



\lipsum[41]

\ce{A + B2O ->[oben][unten] AB2 + O}

\newpage
\part{Selbständigkeitserklärung}
Hiermit erkläre ich, dass ich die vorliegende Hausarbeit selbständig verfasst und keine anderen als die angegebenen Hilfsmittel benutzt habe.\footnote{eine kleine Fußnote.}

Die Stellen der Hausarbeit, die anderen Quellen im Wortlaut oder dem Sinn nach entnommen wurden, sind durch Angaben der Herkunft kenntlich gemacht. Dies gilt auch für Zeichnungen, Skizzen, bildliche Darstellungen sowie für Quellen aus dem Internet.

\vspace{20mm}
\hrule
\vspace{5mm}
\begin{center}
\textbf{Bad Reichenhall}, \today
\end{center}


\newpage
\listoffigures		% Abbildungsverzeichnis erstellen

\newpage
\listoftables		% Tabellenverzeichnis erstellen

\newpage
\addcontentsline{toc}{section}{Literaturverzeichnis}
\small 						%\documentclass[journal = jacsat]{achemso}
							%\setkeys{acs}{articletitle = true}
\bibliographystyle{angewchem} %achemso angewchem 
\bibliography{standard}

\newpage
\appendix
\addcontentsline{toc}{section}{Anhang}
\begin{appendix}
\section*{Anhang}
%\subsection{Anhangsteil 1}
%\lipsum[31]
%\subsection{Anhangsteil 2}
\end{appendix}
\end{document}